\section {The number one half (Algebra)}

A teacher writes $96$ fractions on the board: $\frac{1}{96}, \frac{2}{96}, \frac{3}{96}, \ldots, \frac{95}{96}, \frac{96}{96}$.
A student then plays a game: in every turn he picks any two numbers $a$ and $b$ on the board, erases them, and writes $2ab - a - b + 1$ instead.
After $95$ turns only one number is left. What is it?

\subsection{Solution}

The last number is $\frac{1}{2}$, no matter how the student plays.

Let $f(a,b) = 2ab - a - b + 1$. Then $f(\frac{1}{2}, b) = b - \frac{1}{2} - b + 1 = \frac{1}{2}$, and in the same way $f(a, \frac{1}{2}) = \frac{1}{2}$.
The number $\frac{48}{96} = \frac{1}{2}$ is on the board from the start. Whenever it is picked, the new number is $\frac{1}{2}$ again,
so there is always a $\frac{1}{2}$ on the board, and therefore the last number is $\frac{1}{2}$.

A cleaner way to see what is going on: $2f(a,b) - 1 = 4ab - 2a - 2b + 1 = (2a-1)(2b-1)$.
So if we look at the numbers $2x-1$ for all the numbers $x$ on the board, every turn replaces two of them by their product,
and the product of all of them never changes. At the start one of them is $2 \cdot \frac{1}{2} - 1 = 0$, so the product is $0$,
and at the end the last number $x$ satisfies $2x-1 = 0$.

This is also why the riddle uses an even number like $96$: with an odd number of fractions, $\frac{1}{2}$ would not be on the board.
